\subsection{Cooperative Interception against Maneuvering Targets}
\label{subsec:cases}

Based on the optimal assignment topology, we evaluate three-dimensional cooperative guidance
in two maneuvering scenarios. ART-MAPPO and the capacity-matched PN reference use the same
initial states, target assignment, hit radius, speed limits, load-factor envelope, and other
scenario settings described in Section~\ref{subsec:setup}. For each
method and case, we report the three-dimensional trajectory, its horizontal projection, all
three load commands and their associated attitude/speed responses, time-to-go histories, and
terminal arrival-time differences.

\subsubsection{Case~1: Interception against Evasive Maneuvering Targets}

In Case~1, the attacking UAVs initially execute a constant-load evasive maneuver and transition
to PN guidance after 15\,s. This creates simultaneous horizontal and vertical closing demands
while the defenders must preserve the assignment groups and coordinate their terminal arrival
times.

\emph{1) ART-MAPPO:} Fig.~\ref{fig:case1_mappo_traj_3d} shows the complete spatial
engagement. Starting below the attackers, the defenders climb from the defended region and
close on the eight elevated target trajectories. The horizontal projection in
Fig.~\ref{fig:case1_mappo_traj_xy} preserves the original plan-view interpretation and shows
that all 20 defenders converge on the targets specified by the assignment topology.

\begin{figure}[!ht]
\centering
\includegraphics[width=3.15in]{../new_sim_fig/figures_v10/mappo/mappo_nopn_trajectory_3d.png}
\caption{Three-dimensional engagement trajectories of the defending swarm using ART-MAPPO in
Case~1.}
\label{fig:case1_mappo_traj_3d}
\end{figure}

\begin{figure}[!ht]
\centering
\includegraphics[width=3.15in]{../new_sim_fig/figures_v10/mappo/mappo_nopn_trajectory.png}
\caption{Horizontal projection of the ART-MAPPO engagement trajectories in Case~1.}
\label{fig:case1_mappo_traj_xy}
\end{figure}

The corresponding command and state histories are separated into three physically related
pairs in Figs.~\ref{fig:case1_mappo_yaw}--\ref{fig:case1_mappo_axial}. The horizontal normal
command $n_y$ and heading $\gamma_D$ form the yaw-plane response, while $n_z$ and pitch
$\theta_D$ show the vertical response that is absent from a planar model. ART-MAPPO keeps
$n_y$ inside the $\pm1$ envelope and uses a substantially smaller vertical command for most
of the engagement. The smooth pitch histories confirm gradual altitude capture. Meanwhile,
$n_x$ and $V_D$ show how the policy regulates closing speed without decoupling temporal
coordination from the two normal-load channels.

\begin{figure}[!ht]
\centering
\includegraphics[width=3.1in]{../new_sim_fig/figures_v10/mappo/standalone_duav_legend.png}\\[-2pt]
\subfloat[Normal load $n_y$]{\includegraphics[width=3.2in]{../new_sim_fig/figures_v10/mappo/mappo_nopn_ny.png}}\\[-2pt]
\subfloat[Heading $\gamma_D$]{\includegraphics[width=3.2in]{../new_sim_fig/figures_v10/mappo/mappo_nopn_heading.png}}
\caption{Yaw-plane command and response of ART-MAPPO in Case~1.}
\label{fig:case1_mappo_yaw}
\end{figure}

\begin{figure}[!ht]
\centering
\includegraphics[width=3.1in]{../new_sim_fig/figures_v10/mappo/standalone_duav_legend.png}\\[-2pt]
\subfloat[Vertical load $n_z$]{\includegraphics[width=3.2in]{../new_sim_fig/figures_v10/mappo/mappo_nopn_nz.png}}\\[-2pt]
\subfloat[Pitch $\theta_D$]{\includegraphics[width=3.2in]{../new_sim_fig/figures_v10/mappo/mappo_nopn_pitch.png}}
\caption{Pitch-plane command and response of ART-MAPPO in Case~1.}
\label{fig:case1_mappo_pitch}
\end{figure}

\begin{figure}[!ht]
\centering
\includegraphics[width=3.1in]{../new_sim_fig/figures_v10/mappo/standalone_duav_legend.png}\\[-2pt]
\subfloat[Axial load $n_x$]{\includegraphics[width=3.2in]{../new_sim_fig/figures_v10/mappo/mappo_nopn_nx.png}}\\[-2pt]
\subfloat[Velocity $V_D$]{\includegraphics[width=3.2in]{../new_sim_fig/figures_v10/mappo/mappo_nopn_velocity.png}}
\caption{Axial command and speed response of ART-MAPPO in Case~1.}
\label{fig:case1_mappo_axial}
\end{figure}

The $t_{go}$ histories in Fig.~\ref{fig:case1_mappo_tgo} converge within the eight target
groups despite the three-dimensional engagement geometry. The terminal
summary in Fig.~\ref{fig:case1_mappo_sync} reports a maximum within-group arrival spread of
0.25\,s, below the 0.5-s synchronization criterion.

\begin{figure}[!ht]
\centering
\includegraphics[width=3.2in]{../new_sim_fig/figures_v10/mappo/mappo_nopn_tgo.png}
\caption{Time-to-go ($t_{go}$) profiles of the defending swarm using ART-MAPPO in Case~1.}
\label{fig:case1_mappo_tgo}
\end{figure}

\begin{figure}[!ht]
\centering
\includegraphics[width=3.0in]{../new_sim_fig/figures_v10/mappo/mappo_nopn_time_sync.png}
\caption{Terminal arrival time differences $\Delta t$ within each ART-MAPPO interceptor group
in Case~1.}
\label{fig:case1_mappo_sync}
\end{figure}

\emph{2) PN baseline:} The capacity-matched PN reference is shown in
Figs.~\ref{fig:case1_pn_traj_3d} and~\ref{fig:case1_pn_traj_xy}. PN also completes the selected
engagement, but its purely reactive commands produce visibly different spatial curvature and
control transients. The same command limits are imposed on both methods, so the comparison does
not attribute physically unavailable load authority to the baseline.

\begin{figure}[!ht]
\centering
\includegraphics[width=3.15in]{../new_sim_fig/figures_v10/pn/pn_nopn_trajectory_3d.png}
\caption{Three-dimensional engagement trajectories of the defending swarm using the PN
baseline in Case~1.}
\label{fig:case1_pn_traj_3d}
\end{figure}

\begin{figure}[!ht]
\centering
\includegraphics[width=3.15in]{../new_sim_fig/figures_v10/pn/pn_nopn_trajectory.png}
\caption{Horizontal projection of the PN engagement trajectories in Case~1.}
\label{fig:case1_pn_traj_xy}
\end{figure}

Figures~\ref{fig:case1_pn_yaw}--\ref{fig:case1_pn_axial} give the corresponding yaw-plane,
pitch-plane, and axial responses. The PN commands are clipped to $n_x\in[-0.1,1]$ and
$n_y,n_z\in[-1,1]$, and the command histories show more direct use of the available bounds.
The pitch response remains finite and the speed histories remain within the common vehicle
envelope, providing a physically matched three-dimensional reference.

\begin{figure}[!ht]
\centering
\includegraphics[width=3.1in]{../new_sim_fig/figures_v10/pn/standalone_duav_legend.png}\\[-2pt]
\subfloat[Normal load $n_y$]{\includegraphics[width=3.2in]{../new_sim_fig/figures_v10/pn/pn_nopn_ny.png}}\\[-2pt]
\subfloat[Heading $\gamma_D$]{\includegraphics[width=3.2in]{../new_sim_fig/figures_v10/pn/pn_nopn_heading.png}}
\caption{Yaw-plane command and response of the capacity-matched PN baseline in Case~1.}
\label{fig:case1_pn_yaw}
\end{figure}

\begin{figure}[!ht]
\centering
\includegraphics[width=3.1in]{../new_sim_fig/figures_v10/pn/standalone_duav_legend.png}\\[-2pt]
\subfloat[Vertical load $n_z$]{\includegraphics[width=3.2in]{../new_sim_fig/figures_v10/pn/pn_nopn_nz.png}}\\[-2pt]
\subfloat[Pitch $\theta_D$]{\includegraphics[width=3.2in]{../new_sim_fig/figures_v10/pn/pn_nopn_pitch.png}}
\caption{Pitch-plane command and response of the capacity-matched PN baseline in Case~1.}
\label{fig:case1_pn_pitch}
\end{figure}

\begin{figure}[!ht]
\centering
\includegraphics[width=3.1in]{../new_sim_fig/figures_v10/pn/standalone_duav_legend.png}\\[-2pt]
\subfloat[Axial load $n_x$]{\includegraphics[width=3.2in]{../new_sim_fig/figures_v10/pn/pn_nopn_nx.png}}\\[-2pt]
\subfloat[Velocity $V_D$]{\includegraphics[width=3.2in]{../new_sim_fig/figures_v10/pn/pn_nopn_velocity.png}}
\caption{Axial command and speed response of the capacity-matched PN baseline in Case~1.}
\label{fig:case1_pn_axial}
\end{figure}

The PN $t_{go}$ histories also converge in this selected run, as shown in
Fig.~\ref{fig:case1_pn_tgo}. Its maximum terminal group spread is 0.40\,s
(Fig.~\ref{fig:case1_pn_sync}), compared with 0.25\,s for ART-MAPPO. Both satisfy the 0.5-s
criterion, while the learned policy gives the tighter representative Case~1 synchronization.

\begin{figure}[!ht]
\centering
\includegraphics[width=3.2in]{../new_sim_fig/figures_v10/pn/pn_nopn_tgo.png}
\caption{Time-to-go ($t_{go}$) profiles of the defending swarm using the PN baseline in
Case~1.}
\label{fig:case1_pn_tgo}
\end{figure}

\begin{figure}[!ht]
\centering
\includegraphics[width=3.0in]{../new_sim_fig/figures_v10/pn/pn_nopn_time_sync.png}
\caption{Terminal arrival time differences $\Delta t$ within each PN interceptor group in
Case~1.}
\label{fig:case1_pn_sync}
\end{figure}
\FloatBarrier

\subsubsection{Case~2: Interception against Continuous Weaving Targets}

Case~2 uses a sinusoidal penetration maneuver~\cite{Zarchan1995969}. Its continuous horizontal
and vertical LOS variation provides a more persistent disturbance than the Case~1 transition
and tests whether group timing can be maintained while all three command channels respond.

\emph{1) ART-MAPPO:} The three-dimensional paths in
Fig.~\ref{fig:case2_mappo_traj_3d} show that all eight groups maintain spatial closure against
the weaving targets. The horizontal projection in Fig.~\ref{fig:case2_mappo_traj_xy} exposes
the lateral curvature without hiding the altitude evolution shown in the 3-D view.

\begin{figure}[!ht]
\centering
\includegraphics[width=3.15in]{../new_sim_fig/figures_v10/mappo/mappo_sin_trajectory_3d.png}
\caption{Three-dimensional ART-MAPPO engagement trajectories against continuous weaving
targets in Case~2.}
\label{fig:case2_mappo_traj_3d}
\end{figure}

\begin{figure}[!ht]
\centering
\includegraphics[width=3.15in]{../new_sim_fig/figures_v10/mappo/mappo_sin_trajectory.png}
\caption{Horizontal projection of the ART-MAPPO engagement trajectories in Case~2.}
\label{fig:case2_mappo_traj_xy}
\end{figure}

As shown in Figs.~\ref{fig:case2_mappo_yaw}--\ref{fig:case2_mappo_axial}, ART-MAPPO adapts
$n_y$ to the continuous lateral disturbance while using $n_z$ to regulate vertical closure.
The heading and pitch histories
remain smooth relative to the command variations, and all load channels remain inside the
prescribed envelope. The axial command raises or preserves closing speed without eliminating
the control authority required in the yaw and pitch planes.

\begin{figure}[!ht]
\centering
\includegraphics[width=3.1in]{../new_sim_fig/figures_v10/mappo/standalone_duav_legend.png}\\[-2pt]
\subfloat[Normal load $n_y$]{\includegraphics[width=3.2in]{../new_sim_fig/figures_v10/mappo/mappo_sin_ny.png}}\\[-2pt]
\subfloat[Heading $\gamma_D$]{\includegraphics[width=3.2in]{../new_sim_fig/figures_v10/mappo/mappo_sin_heading.png}}
\caption{Yaw-plane command and response of ART-MAPPO in Case~2.}
\label{fig:case2_mappo_yaw}
\end{figure}

\begin{figure}[!ht]
\centering
\includegraphics[width=3.1in]{../new_sim_fig/figures_v10/mappo/standalone_duav_legend.png}\\[-2pt]
\subfloat[Vertical load $n_z$]{\includegraphics[width=3.2in]{../new_sim_fig/figures_v10/mappo/mappo_sin_nz.png}}\\[-2pt]
\subfloat[Pitch $\theta_D$]{\includegraphics[width=3.2in]{../new_sim_fig/figures_v10/mappo/mappo_sin_pitch.png}}
\caption{Pitch-plane command and response of ART-MAPPO in Case~2.}
\label{fig:case2_mappo_pitch}
\end{figure}

\begin{figure}[!ht]
\centering
\includegraphics[width=3.1in]{../new_sim_fig/figures_v10/mappo/standalone_duav_legend.png}\\[-2pt]
\subfloat[Axial load $n_x$]{\includegraphics[width=3.2in]{../new_sim_fig/figures_v10/mappo/mappo_sin_nx.png}}\\[-2pt]
\subfloat[Velocity $V_D$]{\includegraphics[width=3.2in]{../new_sim_fig/figures_v10/mappo/mappo_sin_velocity.png}}
\caption{Axial command and speed response of ART-MAPPO in Case~2.}
\label{fig:case2_mappo_axial}
\end{figure}

Despite the continuous LOS disturbance, the ART-MAPPO $t_{go}$ curves converge within their
assigned groups (Fig.~\ref{fig:case2_mappo_tgo}). The largest terminal arrival-time difference
is 0.30\,s in Fig.~\ref{fig:case2_mappo_sync}, demonstrating that the timing objective remains
active in the three-dimensional engagement.

\begin{figure}[!ht]
\centering
\includegraphics[width=3.2in]{../new_sim_fig/figures_v10/mappo/mappo_sin_tgo.png}
\caption{Time-to-go ($t_{go}$) profiles of the defending swarm using ART-MAPPO in Case~2.}
\label{fig:case2_mappo_tgo}
\end{figure}

\begin{figure}[!ht]
\centering
\includegraphics[width=3.0in]{../new_sim_fig/figures_v10/mappo/mappo_sin_time_sync.png}
\caption{Terminal arrival time differences $\Delta t$ within each ART-MAPPO interceptor group
in Case~2.}
\label{fig:case2_mappo_sync}
\end{figure}

\emph{2) PN baseline:} Figures~\ref{fig:case2_pn_traj_3d}
and~\ref{fig:case2_pn_traj_xy} show that the constrained PN baseline also reaches all assigned
targets in the selected Case~2 engagement. The spatial paths exhibit the reactive response to
the continuously varying LOS, while the plan view and 3-D view together distinguish lateral
curvature from vertical closure.

\begin{figure}[!ht]
\centering
\includegraphics[width=3.15in]{../new_sim_fig/figures_v10/pn/pn_sin_trajectory_3d.png}
\caption{Three-dimensional engagement trajectories of the defending swarm using the PN
baseline in Case~2.}
\label{fig:case2_pn_traj_3d}
\end{figure}

\begin{figure}[!ht]
\centering
\includegraphics[width=3.15in]{../new_sim_fig/figures_v10/pn/pn_sin_trajectory.png}
\caption{Horizontal projection of the PN engagement trajectories in Case~2.}
\label{fig:case2_pn_traj_xy}
\end{figure}

The PN profiles in Figs.~\ref{fig:case2_pn_yaw}--\ref{fig:case2_pn_axial} show sustained
reactive changes in the horizontal channel together with the vertical and axial responses.
Since all commands are
clipped to the same envelope as ART-MAPPO, excursions that reach the plotted $\pm1$ lines
represent saturation of available authority rather than an unrealizable command.

\begin{figure}[!ht]
\centering
\includegraphics[width=3.1in]{../new_sim_fig/figures_v10/pn/standalone_duav_legend.png}\\[-2pt]
\subfloat[Normal load $n_y$]{\includegraphics[width=3.2in]{../new_sim_fig/figures_v10/pn/pn_sin_ny.png}}\\[-2pt]
\subfloat[Heading $\gamma_D$]{\includegraphics[width=3.2in]{../new_sim_fig/figures_v10/pn/pn_sin_heading.png}}
\caption{Yaw-plane command and response of the capacity-matched PN baseline in Case~2.}
\label{fig:case2_pn_yaw}
\end{figure}

\begin{figure}[!ht]
\centering
\includegraphics[width=3.1in]{../new_sim_fig/figures_v10/pn/standalone_duav_legend.png}\\[-2pt]
\subfloat[Vertical load $n_z$]{\includegraphics[width=3.2in]{../new_sim_fig/figures_v10/pn/pn_sin_nz.png}}\\[-2pt]
\subfloat[Pitch $\theta_D$]{\includegraphics[width=3.2in]{../new_sim_fig/figures_v10/pn/pn_sin_pitch.png}}
\caption{Pitch-plane command and response of the capacity-matched PN baseline in Case~2.}
\label{fig:case2_pn_pitch}
\end{figure}

\begin{figure}[!ht]
\centering
\includegraphics[width=3.1in]{../new_sim_fig/figures_v10/pn/standalone_duav_legend.png}\\[-2pt]
\subfloat[Axial load $n_x$]{\includegraphics[width=3.2in]{../new_sim_fig/figures_v10/pn/pn_sin_nx.png}}\\[-2pt]
\subfloat[Velocity $V_D$]{\includegraphics[width=3.2in]{../new_sim_fig/figures_v10/pn/pn_sin_velocity.png}}
\caption{Axial command and speed response of the capacity-matched PN baseline in Case~2.}
\label{fig:case2_pn_axial}
\end{figure}

Figure~\ref{fig:case2_pn_tgo} shows that the selected PN engagement also achieves groupwise
time-to-go convergence. Its maximum arrival spread is 0.35\,s
(Fig.~\ref{fig:case2_pn_sync}), whereas ART-MAPPO gives 0.30\,s. These representative runs
therefore show successful three-dimensional interception for both methods and slightly tighter
terminal timing for ART-MAPPO; the Monte Carlo analysis below remains the basis for broader
reliability comparisons.

\begin{figure}[!ht]
\centering
\includegraphics[width=3.2in]{../new_sim_fig/figures_v10/pn/pn_sin_tgo.png}
\caption{Time-to-go ($t_{go}$) profiles of the defending swarm using the PN baseline in
Case~2.}
\label{fig:case2_pn_tgo}
\end{figure}

\begin{figure}[!ht]
\centering
\includegraphics[width=3.0in]{../new_sim_fig/figures_v10/pn/pn_sin_time_sync.png}
\caption{Terminal arrival time differences $\Delta t$ within each PN interceptor group in
Case~2.}
\label{fig:case2_pn_sync}
\end{figure}
\FloatBarrier
